% Part: first-order-logic
% Chapter: axiomatic-deduction
% Section: provability-quantifiers

% verification of properties of provability needed for maximally
% consistent sets in the completeness chapter.

\documentclass[../../../include/open-logic-section]{subfiles}

\begin{document}

\olfileid{fol}{axd}{qpr}

\olsection{\usetoken{S}{derivability} మరియు పరిమాణీకరణికాలు}

\begin{explain}
  ఈ విభాగంలో స్థాపించే~$\Proves$ గురించిన వాస్తవాలను స్వీకృతాధారిత
  నిగమనాలు ఇవ్వాలని కూడా సంపూర్ణతా సిద్ధాంతం కోరుతుంది.
\end{explain}

\begin{thm}
\ollabel{thm:strong-generalization} $c$ అనేది $\Gamma$లో గానీ
$!A(x)$లో గానీ సంభవించని ఒక \tetoken{స్థిరసంకేతం}{!!a{constant}} అయ్యి, $\Gamma \Proves
!A(c)$ అయితే, $\Gamma \Proves \lforall[x][!A(x)]$.
\end{thm}

\begin{proof}
నిగమన సిద్ధాంతం వల్ల $\Gamma \Proves \ltrue \lif !A(c)$. $c$ అనేది
$\Gamma$లో గానీ~$\top$లో గానీ సంభవించదు కనుక, $\Gamma \Proves
\ltrue \lif \lforall[x][!A(x)]$. $\Gamma \Proves \ltrue$ అనేది
\texttt{ax:ltrue} స్వీకృత రూపం; కాబట్టి మోడస్ పోనెన్స్ వల్ల $\Gamma
\Proves \lforall[x][!A(x)]$.
\sourcecorrection{OLTEAXD-009}{ఆంగ్ల మూలం చివరి దశను ``నిగమన సిద్ధాంతం
మళ్లీ'' అని సమర్థించింది; అయితే $\Gamma \Proves \ltrue \lif
\lforall[x][!A(x)]$ నుంచి కావలసిన ఫలితం రావడానికి
$\Gamma \Proves \ltrue$ అనే స్వీకృత రూపంతో మోడస్ పోనెన్స్ అవసరం.
ఆ రెండు దశలను స్పష్టంగా ఇచ్చాం.}
\end{proof}

\begin{prop}
\ollabel{prop:provability-quantifiers}
\begin{tagenumerate}{prvEx,prvAll}
\tagitem{prvEx}{$!A(t) \Proves \lexists[x][!A(x)]$.}{}

\tagitem{prvAll}{$\lforall[x][!A(x)] \Proves !A(t)$.}{}
\end{tagenumerate}
\end{prop}

\begin{proof}
\begin{tagenumerate}{prvEx,prvAll}
\tagitem{prvEx}{\olref[qua]{ax:q2}, నిగమన సిద్ధాంతం ద్వారా.}{}
\tagitem{prvAll}{\olref[qua]{ax:q1}, నిగమన సిద్ధాంతం ద్వారా.}{}
\end{tagenumerate}
\end{proof}

\end{document}
